\scanpage{34}
% archive sheet number: 33
\begingroup
\makeatletter\tagsleft@true\makeatother
% loc: opening, (13)
Posons
\[
\Sigma_{12}=\Sigma_1^-\cap\Sigma_2^-,\qquad E_1=\Sigma_1^+\cap\Sigma_1^-,\quad E_2=\Sigma_1^+\cap\Sigma_2^-\tag*{(13)}
\]
alors on a le diagramme d'inclusion de parties fermées de $X$
% loc: diagram (14), lower slant
\usetikzlibrary{decorations.pathmorphing}
\diagram{\[
\begin{tikzcd}[ampersand replacement=\&,row sep=1.1em,column sep=.4em]
\& \& X \& \& \\
\& \Sigma_2^- \arrow[ur] \& \& \Sigma_1^- \arrow[ul] \& \\
\Sigma_1^+ \arrow[ur] \& \& \Sigma_{12} \arrow[ul]\arrow[ur] \& \& \Sigma_2^+ \arrow[ul]\\
\& E_1 \arrow[ul]\arrow[ur] \& \& E_2 \arrow[ul]\arrow[ur] \& \\
\& \& \deleted{\unsure{\Omega|2|\varnothing}}\arrow[ul]\arrow[ul,no head,decorate,decoration={zigzag,segment length=3.4pt,amplitude=1.6pt}] \& \&
\end{tikzcd}\tag*{(14)}
\]}
% loc: right marginal NB alongside (14)
\margin{NB $\Sigma_1^+\cap\Sigma_2^+=\varnothing$ donc $E_1\cap\Sigma_2^+=\varnothing$, $E_2\cap\Sigma_1^-=\varnothing$ i.e. $E_1\subset\Sigma_2^{\unsure{\circ|o|0}-}$, $E_2\subset\Sigma_1^{\circ-}$ et $E_1\cap E_2\unsure{.|-}=\varnothing$}
% loc: line under (14), cancelled starter and next display
où \unsure{les|le} trois carrés sont cartésiens et cocartésiens
\[
\deleted{\unsure{E_1|E_+}}
\]
i.e.
\[
\left|\begin{aligned}
&E_1=\Sigma_1^+\cap\Sigma_{12},\quad\Sigma_2^-=\Sigma_1^+\cup\Sigma_{12}\\
&\deleted{\unsure{E_1|E}}
\end{aligned}\right.
\]
et les deux relations analogues correspondant aux deux autres carrés
% loc: (15), writes over first row and lower X
\[
\left\{\begin{aligned}
&E_2=\Sigma_{\deleted{\unsure{2|1}}12}^{\deleted{\unsure{+|-}}}\cap\Sigma_{\deleted{\unsure{12|1}}2}^+,\quad \Sigma_1^-=\Sigma_{\unsure{22|12}}^{\deleted{\unsure{+|-}}}\cup\Sigma_2^+\\
&\Sigma_{12}\overset{(13)}{\unsure{\equiv|=}}\Sigma_1^-\cap\Sigma_{\unsure{2|1}}^-,\quad \deleted{\unsure{\Sigma|{-}}}X=\Sigma_2^-\cup\Sigma_1^-.
\end{aligned}\right.\tag*{(15)}
\]

% loc: (16), three lines
Notons que
\[
\left\{\begin{aligned}
&\text{a)}\quad\Sigma_1^+\cap\Sigma_2^-=\Sigma_1^+\ne\varnothing\\
&\text{b)}\quad\Sigma_2^+\cap\Sigma_1^-=\Sigma_2^+\ne\varnothing\\
&\text{c)}\quad\Sigma_1^-\cap\Sigma_2^-=\Sigma_{12}\ne\varnothing,
\end{aligned}\right.\tag*{(16)}
\]

% loc: after (16), interline addition
\inserted{\unsure{Prouvons|Prenons}} \unsure{La|la} dernière relation (les deux premières étant évidentes).
% loc: last paragraph; abandoned inline and display formulas
Si on a\unsure{\deleted{vait}|vait|\deleted{ura}} $\Sigma_{12}=\varnothing$, comme $E_1,E_2\subset\Sigma_{12}$, on aurait $E_1=\varnothing$, $E_2=\varnothing$, de sorte que $\pi_1$ et $\pi_2$ sont deux \uline{partitions} de $X$ en des parties non vides, i.e. $\deleted{\Sigma_{\unsure{1|2}}^- = X\setminus{\unsure{\Sigma|{-}}}^-}$ on \unsure{aurait|avait}
\[
\deleted{\unsure{\Sigma_1^+|\Sigma_1^-}}\quad\Sigma_1^-=X\setminus\Sigma_1^+,\quad\Sigma_2^-=X\unsure{\setminus|-}\Sigma_{\unsure{2|1}}^+
\]
\endgroup
